\documentclass[12pt]{book}
\usepackage[utf8]{inputenc}
\usepackage[T1]{fontenc}
\usepackage{amsmath,amssymb,amsfonts}
\usepackage{mathrsfs}
\usepackage{amsthm}

% Calligraphic for categories and sheaves
\newcommand{\cat}[1]{\mathcal{#1}}
\newcommand{\sheaf}[1]{\mathcal{#1}}

% Standard math operators
\DeclareMathOperator{\Hom}{Hom}
\DeclareMathOperator{\Ext}{Ext}
\DeclareMathOperator{\Tor}{Tor}
\DeclareMathOperator{\id}{id}
\DeclareMathOperator{\varinjlim}{\varinjlim}
\DeclareMathOperator{\RHom}{\mathbf{R}\mathcal{H}\textit{om}}
\DeclareMathOperator{\RGamma}{\mathbf{R}\Gamma}

\begin{document}

\setcounter{page}{276}
\section{Local Duality.}

Throughout this section we let \(A\) be a noetherian local ring, with maximal ideal \(\mathfrak{m}\), and residue field \(k\), and let \(R^{\cdot}\) be a dualizing complex on \(X = \operatorname{Spec}A\) (cf. 62). Recall
Proposition 3.4 which gives a necessary and sufficient condition for a complex \(R^{\cdot} \in D_{C}^{+}(A)\) to be dualizing.

\textbf{Definition.}
We say that the dualizing complex \(R^{\cdot}\) is normalized if the integer \(d\) of Proposition 3.4 is zero.

\textbf{Remark.}
Since the translate of a dualizing complex is again one, we can normalize by translation.

\begin{proposition}
If \(R^{\cdot}\) is normalized, then \(\RGamma_{x}(R^{\cdot})\), where \(x\) is the closed point of \(\operatorname{Spec}A\), is isomorphic in the derived category to an injective hull \(I\) of \(k\).
\end{proposition}

\begin{proof}
The cohomology of \(\RGamma_{x}(R^{\cdot})\) is
\[H_{x}^{i}\left(R^{\cdot}\right) \cong \varinjlim_{n} Ext^{i}\left(A / \mathfrak{m}^{n}, R^{\cdot}\right)\]
by Corollary 4.2 above. On the other hand, we saw from the proof of Proposition 3.4 that \(Ext^{i}(M, R^{\cdot})=0\) for \(i \neq0\) and all \(M\) of finite length. Furthermore, by Proposition 5.2,the functor \(Ext^{\circ}(\cdot, R^{\cdot})\), is a dualizing functor, and corresponds to

\setcounter{page}{277}
the injective hull of \(k\) given by
\[I=\varinjlim_{n} Ext^{\circ}\left(A / \mathfrak{m}^{n}, R^{\cdot}\right) .\]

So we see that \(H_{x}^{i}(R^{\cdot})=0\) for \(i \neq0\),and \(H_{x}^{\circ}(R^{\cdot})=I\).
Therefore \(\RGamma_{x}(R^{\cdot}) \cong I\) in the derived category [I.4.3].
\end{proof}

Now let \(M\) be an \(A\)-module.
Then there is a natural homomorphism
\[\Gamma_{x}(M) \to \Hom\left(\Hom\left(M, R^{\cdot}\right), \Gamma_{x}\left(R^{\cdot}\right)\right) .\]

This gives rise to a morphism of functors on \(D^{+}(A)\),
\[\RGamma_{x}\left(M^{\cdot}\right) \to \RHom\left(\RHom\left(M^{\cdot}, R^{\cdot}\right), \RGamma_{x}\left(R^{\cdot}\right)\right),\]
since we can take \(R^{\cdot}\) to be an injective complex, in which case \(\RGamma_{x}(R^{\cdot})=\Gamma_{x}(R^{\cdot})\) is also injective. Indeed, whenever \(J\) is an injective \(A\)-module, and \(Y\) a closed subset of \(\operatorname{Spec}A\), then \(\Gamma_{Y}(J)\) is injective (cf. [II] for the structure of injective \(A\)-modules). We now assume (without loss of generality) that \(R^{\cdot}\) is normalized, and so, by the Proposition, have a morphism of functors

\setcounter{page}{278}
\[\theta: \RGamma_{x}\left(M^{\cdot}\right) \to \Hom\left(\RHom\left(M^{\cdot}, R^{\cdot}\right), I\right) .\]

(Following the usual conventions, we do not write \(\mathbf{R}\) before a functor which is already exact.)

\begin{theorem}[Local Duality]
Let \(A\) be a local noetherian ring, let \(R^{\cdot}\) be a normalized dualizing complex, let \(I\) be the corresponding injective hull of \(k\), and let \(M^{\cdot} \in D_{C}^{+}(A)\) (i.e., the cohomology modules of \(M^{\cdot}\) are of finite type). Then \(\theta\) as defined above is an isomorphism.
\end{theorem}

\begin{corollary}
With \(A\), \(R^{\cdot}\), \(I\) as above, let \(M\) be a module of finite type. Then the homomorphisms
\[\theta^{i}: H_{x}^{i}(M) \to \Hom\left(Ext^{-i}\left(M, R^{\cdot}\right), I\right)\]
induced by \(\theta\) are isomorphisms.
\end{corollary}

\begin{proof}
Using the Lemma on Way-out Functors [I.?.l] one sees that the Corollary is equivalent to the theorem. To prove the Corollary, note that \(\theta^{i}\) is an isomorphism for \(M=k\), because \(H_{x}^{i}(k)=0\) for \(i \neq0\),and \(H_{x}^{\circ}(k)=k\); note also that for all \(M\) of finite type, both sides are modules with support at \(x\). Thus we are reduced to proving the following

\setcounter{page}{279}
\begin{lemma}
Let \(\theta^{i}: S^{i} \to T^{i}\) be a morphism of covariant cohomological functors on the category of modules over a noetherian ring \(A\). Assume
\begin{enumerate}
\item[(i)] for every maximal ideal \(\mathfrak{m} \subseteq A\), \(\theta^{i}(A / \mathfrak{m})\) is an isomorphism,
\item[(ii)] for every non-maximal prime ideal \(\mathfrak{p} \subseteq A\), \(S^{i}(A / \mathfrak{p})\) and \(T^{i}(A / \mathfrak{p})\) have support \(< \operatorname{Supp}(A/\mathfrak{p})\).
\end{enumerate}
Then \(\theta^{i}(M)\) is an isomorphism for every \(A\)-module \(M\) of finite type.
\end{lemma}

\begin{proof}
If \(M\) is any module of finite type, then admits a finite filtration each of whose quotients is of the form \(A/P\) with \(P\) prime. By assumption, \(\theta^{i}(A / \mathfrak{m})\) is an isomorphism for every maximal ideal \(\mathfrak{m}\). Thus by the 5-lemma we reduce to the case \(M=A / \mathfrak{p}\) with \(\mathfrak{p}\) not maximal. By noetherian induction, we may assume that \(\theta^{i}(M')\) is an isomorphism for every \(M'\) of finite type with support \(< \operatorname{Supp}(A/\mathfrak{p})\).

Now for each \(f \in M=A / \mathfrak{p}\), \(f \neq0\), we consider the exact sequence
\[0 \to M \stackrel{f}{\to} M \to M / fM \to 0,\]
and apply the functors \(S^{i}\), \(T^{i}\) to it, splitting it as follows:

\setcounter{page}{280}
\[\cdots \to S^{i}(M) \stackrel{f}{\to} S^{i}(M) \to \cdots\]
and
\[\cdots \to S^{i-1}(M) \to S^{i-1}(M) \to S^{i-1}(M / f M) \to K_{f} \to 0\]

Now by our induction hypothesis, is an isomorphism, so \(\alpha\) is surjective. Since \(\operatorname{Supp}S^{i}(M)<\operatorname{supp} M\), each element of \(S^{i}(M)\) is annihilated by some non-zero \(f \in M\).
Thus \(S^{i}(M)=\bigcup K_{f}\) as \(f\) ranges over the non-zero elements of \(M\). Similarly \(T^{i}(M)=\bigcup L_{f}\). Therefore \(\beta\) is surjective.

This is true for each \(i\), so also \(\gamma\) is surjective, But that implies \(\alpha\) is injective, so \(\alpha\) is an isomorphism. Then, as above, is also an isomorphism. q.e.d.
\end{proof}

\begin{corollary}
For \(M\) of finite type, the modules \(Ext^{i}\left(M, R^{\cdot}\right)\) are of co-finite type (see Proposition 5.3), and we have a functorial isomorphism
\[Ext^{i}\left(M, R^{\cdot}\right)^{\wedge} \to \Hom\left(H_{x}^{-i}(M), I\right)\]
of \(\hat{A}\)-modules.
\end{corollary}

\begin{proof}
Apply the functor \(\Hom(\cdot, I)\) to the isomorphism \(\theta^{i}\), and use the fact that for any \(A\)-module of finite type \(M\), this functor applied twice gives the completion \(\hat{M}\). [Ic p. 61].
\end{proof}

\end{document}